\clearpage
\begingroup
\fancyhead[L]{\small NACHLASS UND BRIEFWECHSEL}
\renewcommand{\headrulewidth}{0pt}
\par\medskip
\Needspace{8\baselineskip}
\pdfbookmark[2]{[5] Nicolai an Gauss, 6. August 1814}{pallas-VIII-letter5}
\begin{center}
[5.]\par\medskip
\textsc{Nicolai} an \textsc{Gauss}, Seeberg, 6. August 1814.
\end{center}
\par\noindent\hspace*{\parindent}\{.\ .\ .\ .\ .\quad Was die Pallasstörungen durch Saturn anbetrifft, so habe ich
zwar noch keine Prüfung durch Berechnung einiger speciellen Störungen vorgenommen
(indem ich seit meinem Wiederhiersein mich fast gar nicht mit
numerischen Rechnungen, sondern dafür mit einigen theoretischen Sachen beschäftigt
habe), dagegen aber den ganzen Gang der Rechnung noch einmal
sorgfältig übersehen, allein ich habe keine Unrichtigkeiten darin bemerkt, wie
ich wohl bei der Sorgfalt, mit der ich bei der ganzen Rechnung verfahren
zu sein mir bewusst war, hätte vorher sagen wollen. Auch mit der Annahme
des Zeichens der Kräfte hat es seine völlige Richtigkeit; indem ich nachher,
ich weiss nicht mehr aus was für einem Grunde, ebenfalls von der Annahme
\(\dfrac{x'-x}{\rho^3}-\dfrac{x'}{r^{\prime3}}\) ausgegangen bin; und Sie werden hoffentlich bei allen Elementen
Ihre Coefficienten mit den meinigen übereinstimmend finden. Noch einen
Beruhigungsgrund für die Richtigkeit meiner Gleichungen schöpfe ich daher,
dass die grössern von ihnen ziemlich gut mit meiner frühern Rechnung, die
\SourceLeaf{591}
doch nach einer ganz verschiedenen Methode geführt war, übereinstimmen.
Indess will ich nächstens ein paar specielle Störungen von einer Opposition
zur andern berechnen, und ich hoffe auch hier eine schöne Übereinstimmung
zu erhalten. Sollte diese nicht erfolgen, so wäre ich dennoch lieber geneigt,
die Schuld auf die numerische Berechnung der Gleichungen, als auf diese
selbst, überzutragen. Wer weiss, ob nicht auch die Marsstörungen die Übereinstimmung
der bisher beobachteten Pallasoppositionen, wenn auch nur etwas,
verschlimmern werden?\quad .\ .\ .\ .\ .\}

\par\medskip
\Needspace{8\baselineskip}
\pdfbookmark[2]{[6] Nicolai an Gauss, 18. März und 17. Julius 1815}{pallas-VIII-letters6}
\begin{center}
[6.]\par\medskip
\textsc{Nicolai} an \textsc{Gauss}, Seeberg, 18. März 1815.
\end{center}
\par\noindent\hspace*{\parindent}\{.\ .\ .\ .\ .\quad Ihren Auftrag, die beiden Theile der Epochenstörung für Pallas
zusammenzuschmelzen, werde ich nächstens in Erfüllung bringen und Ihnen
die Formeln zusenden. Für die Störung durch Jupiter muss dies eine ganz
artige Arbeit gewesen sein.\quad .\ .\ .\ .\ .\}

\par\medskip
\Needspace{6\baselineskip}
\begin{center}
\textsc{Nicolai} an \textsc{Gauss}, Seeberg, 17. Julius 1815.
\end{center}
\par\noindent\hspace*{\parindent}\{Gestern Nachmittag habe ich mich mit der Vereinigung der beiden
Theile der Störung der Epoche der Pallas durch Saturn beschäftigt, und ich
eile, Ihnen die gefundenen Gleichungen mitzutheilen. Diejenigen Gleichungen
im ersten Theile der Epoche, zu denen sich im zweiten keine correspondirenden
fanden, waren sämmtlich unter \(0\overset{\prime\prime}{,}1\), und ich habe sie also unter den
Endgleichungen für die Epoche nicht mit aufgeführt. Bei dieser neuen Form
der Epochenstörung wird nun aber manche kleine Gleichung \textls[100]{über} \(0\overset{\prime\prime}{,}1\) fehlen,
deren secrete Theile gleichwohl \textls[100]{unter} \(0\overset{\prime\prime}{,}1\) sind. Besonders wird dies bei
Ihren Jupitersstörungen häufig der Fall sein. --- Im Laufe dieser Arbeit habe
ich bemerkt, dass, wenn man den ersten Theil der Epochenstörung durch
Addition von \(90^0\) zu dem constanten Winkel ebenfalls wie den zweiten auf
die Form \(a\sin(\varphi+A)\) bringt, alsdann, besonders bei den grössern Gleichungen,
jene constanten Winkel nicht sehr von einander verschieden sind. Ist diese
Erscheinung blosser Zufall, oder hat sie irgend einen theoretischen Grund?
--- und findet etwas Ähnliches auch bei den Jupitersstörungen Statt?\quad .\ .\ .\ .\ .

\SourceLeaf{592}
\par\medskip
\Needspace{13\baselineskip}
\begin{center}
Vollständige Störung der Epoche der Pallas durch Saturn.\par
Jährliche Änderung \(=-1\overset{\prime\prime}{,}31973\)
\end{center}
\begingroup
\small
\newcommand{\VIIIpallas}{\text{\AstroSymbols ⚴}}
\newcommand{\VIIIsaturn}{\text{\AstroSymbols ♄}}
\setlength{\arraycolsep}{2pt}
\renewcommand{\arraystretch}{1.25}
\[
\begin{array}{r@{\,}l@{\qquad}r@{\,}l}
0\overset{\prime\prime}{,}83&\sin(\VIIIpallas+175^0\ 20')&+0\overset{\prime\prime}{,}50&\sin(3\VIIIpallas-2\VIIIsaturn+63^0\ 14')\\
+0{,}53&\sin(2\VIIIpallas+182\ 41)&+0{,}30&\sin(3\VIIIsaturn+207\ 41)\\
+1{,}78&\sin(\VIIIsaturn+140\ 25)&+2{,}41&\sin(\VIIIpallas-3\VIIIsaturn+169\ 38)\\
+0{,}47&\sin(\VIIIpallas+\VIIIsaturn+180\ 3)&+0{,}71&\sin(2\VIIIpallas-3\VIIIsaturn+162\ 17)\\
+0{,}61&\sin(\VIIIpallas-\VIIIsaturn+179\ 37)&+0{,}66&\sin(3\VIIIpallas-3\VIIIsaturn+172\ 34)\\
+0{,}18&\sin(2\VIIIpallas-\VIIIsaturn+161\ 42)&+0{,}46&\sin(\VIIIpallas-4\VIIIsaturn+187\ 14)\\
+2{,}01&\sin(2\VIIIsaturn+221\ 24)&+0{,}24&\sin(2\VIIIpallas-4\VIIIsaturn+164\ 16)\\
+7{,}26&\sin(\VIIIpallas-2\VIIIsaturn+124\ 36)&+0{,}15&\sin(3\VIIIpallas-4\VIIIsaturn+161\ 42).\}\\
+3{,}69&\sin(2\VIIIpallas-2\VIIIsaturn+178\ 30)&&
\end{array}
\]
\endgroup
\begin{center}
\rule{0.3\linewidth}{0.4pt}
\end{center}
\clearpage
\endgroup
